\topictitle{\tr{Inducția matematică}{Mathematical induction}{Vollständige Induktion}}

\[
P(n):\ \dots
\]
\begin{enumerate}[label=\arabic*)]
\item \textbf{\tr{Verificare}{Base case}{Induktionsanfang}:}\ \tr{verificăm că $P(n_0)$ este „Adevărat''.}{check that $P(n_0)$ is ``True''.}{man prüft, dass $P(n_0)$ „wahr`` ist.}
\item \textbf{\tr{Demonstrație}{Induction step}{Induktionsschritt}:}\ \tr{presupunem că $P(n)$ este „Adevărat'' și demonstrăm că $P(n+1)$ este „Adevărat''.}
  {assume that $P(n)$ is ``True'' and prove that $P(n+1)$ is ``True''.}
  {man nimmt an, dass $P(n)$ „wahr`` ist, und zeigt, dass $P(n+1)$ „wahr`` ist.}
\item \textbf{\tr{Concluzie}{Conclusion}{Schluss}:}\ $P(n)$ \tr{este adevărată pentru orice $n\ge n_0$.}{is true for all $n\ge n_0$.}{gilt für alle $n\ge n_0$.}
\end{enumerate}
\[
P(n):\ \dots\qquad P(n+1):\ \dots
\]
\hrule\vspace{4pt}
\tr{Pentru a trece de la $P(n)$ la $P(n+1)$ se compară cele două expresii prin scădere ($-$) sau împărțire ($\div$).}
    {To get from $P(n)$ to $P(n+1)$, compare the two expressions by subtraction ($-$) or division ($\div$).}
    {Beim Übergang von $P(n)$ zu $P(n+1)$ vergleicht man die beiden Ausdrücke durch Subtraktion ($-$) oder Division ($\div$).}
